\documentclass[12pt,a4paper]{article}

% --- Packages ---
\usepackage[utf8]{inputenc}
\usepackage[T1]{fontenc}
\usepackage[french]{babel}
\usepackage{geometry}
\usepackage{enumitem}
\usepackage{array}
\usepackage{fancyhdr}
\usepackage{amsmath, amssymb}

% --- Mise en page ---
\geometry{margin=2.5cm}

% --- Pied de page ---
\pagestyle{fancy}
\fancyhf{}
\fancyfoot[R]{\it{G.M.}}
\renewcommand{\headrulewidth}{0pt}
\renewcommand{\footrulewidth}{0pt}

\setlength{\parindent}{0pt}

\begin{document}

\begin{center}
    \Large \textbf{--Automatismes - QCM--} \\[5pt]
    \normalsize Aucune justification n'est demandée et une seule réponse est possible par question.
\end{center}

\vspace{0.7cm}

\textbf{Question 1 :}
\vspace{5pt}

Dans une classe de 40 élèves, 8 sont gauchers. La proportion d'élèves gauchers est :
\vspace{0.3cm}

\renewcommand{\arraystretch}{3}
\begin{tabular}{|m{3.4cm}|m{3.4cm}|m{3.4cm}|m{3.4cm}|}
\hline
\centering \vspace{-0.3cm} \textbf{A.} 0,3 & 
\centering \vspace{-0.3cm} \textbf{B.} 0,25 & 
\centering \vspace{-0.3cm} \textbf{C.} $\dfrac{1}{5}$ & 
\centering \vspace{-0.3cm} \textbf{D.} 0,35 \tabularnewline
\hline
\end{tabular}

\vspace{0.7cm}

\textbf{Question 2 :}
\vspace{5pt}

Un article coûte 240 €. Le vendeur propose une réduction de 15\%. Combien d'euros représente la réduction ?
\vspace{0.3cm}

\renewcommand{\arraystretch}{3}
\begin{tabular}{|m{3.4cm}|m{3.4cm}|m{3.4cm}|m{3.4cm}|}
\hline
\centering \vspace{-0.3cm} \textbf{A.} 30 € & 
\centering \vspace{-0.3cm} \textbf{B.} 24 € & 
\centering \vspace{-0.3cm} \textbf{C.} 42 € & 
\centering \vspace{-0.3cm} \textbf{D.} 36 € \tabularnewline
\hline
\end{tabular}

\vspace{0.7cm}

\textbf{Question 3 :}
\vspace{5pt}

Dans un collège, 18\% des élèves sont en classe de 3ème, ce qui représente 54 élèves. L'effectif total du collège est :
\vspace{0.3cm}

\renewcommand{\arraystretch}{3}
\begin{tabular}{|m{3.4cm}|m{3.4cm}|m{3.4cm}|m{3.4cm}|}
\hline
\centering \vspace{-0.3cm} \textbf{A.} 250 & 
\centering \vspace{-0.3cm} \textbf{B.} 300 & 
\centering \vspace{-0.3cm} \textbf{C.} 350 & 
\centering \vspace{-0.3cm} \textbf{D.} 280 \tabularnewline
\hline
\end{tabular}

\vspace{0.7cm}

\textbf{Question 4 :}
\vspace{5pt}

Un prix augmente de 12\%. Le coefficient multiplicateur associé à cette augmentation est :

\vspace{0.3cm}

\renewcommand{\arraystretch}{3}
\begin{tabular}{|m{3.4cm}|m{3.4cm}|m{3.4cm}|m{3.4cm}|}
\hline
\centering \vspace{-0.3cm} \textbf{A.} 0,12 & 
\centering \vspace{-0.3cm} \textbf{B.} 1,12 & 
\centering \vspace{-0.3cm} \textbf{C.} 0,88 & 
\centering \vspace{-0.3cm} \textbf{D.} 1,2 \tabularnewline
\hline
\end{tabular}

\vspace{0.7cm}

\textbf{Question 5 :}
\vspace{5pt}

Un article coûte 110 €. Après une augmentation de 10\%, son nouveau prix est :
\vspace{0.3cm}

\renewcommand{\arraystretch}{3}
\begin{tabular}{|m{3.4cm}|m{3.4cm}|m{3.4cm}|m{3.4cm}|}
\hline
\centering \vspace{-0.3cm} \textbf{A.} 120 € & 
\centering \vspace{-0.3cm} \textbf{B.} 117 € & 
\centering \vspace{-0.3cm} \textbf{C.} 119 € & 
\centering \vspace{-0.3cm} \textbf{D.} 121 € \tabularnewline
\hline
\end{tabular}

\vspace{0.7cm}

\textbf{Question 6 :}
\vspace{5pt}

Un article coûte 180 €. Après une réduction de 10\%, son nouveau prix est :
\vspace{0.3cm}

\renewcommand{\arraystretch}{3}
\begin{tabular}{|m{3.4cm}|m{3.4cm}|m{3.4cm}|m{3.4cm}|}
\hline
\centering \vspace{-0.3cm} \textbf{A.} 160 € & 
\centering \vspace{-0.3cm} \textbf{B.} 172 € & 
\centering \vspace{-0.3cm} \textbf{C.} 162 € & 
\centering \vspace{-0.3cm} \textbf{D.} 168 € \tabularnewline
\hline
\end{tabular}

\vspace{0.7cm}

\textbf{Question 7 :}
\vspace{5pt}

Après une augmentation de 14\%, un article coûte 114 €. Son prix initial était :
\vspace{0.3cm}

\renewcommand{\arraystretch}{3}
\begin{tabular}{|m{3.4cm}|m{3.4cm}|m{3.4cm}|m{3.4cm}|}
\hline
\centering \vspace{-0.3cm} \textbf{A.} 100 € & 
\centering \vspace{-0.3cm} \textbf{B.} 98 € & 
\centering \vspace{-0.3cm} \textbf{C.} 102 € & 
\centering \vspace{-0.3cm} \textbf{D.} 105 € \tabularnewline
\hline
\end{tabular}

\vspace{0.7cm}

\textbf{Question 8 :}
\vspace{5pt}

Après une réduction de 25\%, un article coûte 75 €. Son prix initial était :
\vspace{0.3cm}

\renewcommand{\arraystretch}{3}
\begin{tabular}{|m{3.4cm}|m{3.4cm}|m{3.4cm}|m{3.4cm}|}
\hline
\centering \vspace{-0.3cm} \textbf{A.} 100 € & 
\centering \vspace{-0.3cm} \textbf{B.} 90 € & 
\centering \vspace{-0.3cm} \textbf{C.} 110 € & 
\centering \vspace{-0.3cm} \textbf{D.} 95 € \tabularnewline
\hline
\end{tabular}

\vspace{0.7cm}

\textbf{Question 9 :}
\vspace{5pt}

Un prix passe de 80 € à 92 €. Le taux d'évolution, exprimé en pourcentage, est :
\vspace{0.3cm}

\renewcommand{\arraystretch}{3}
\begin{tabular}{|m{3.4cm}|m{3.4cm}|m{3.4cm}|m{3.4cm}|}
\hline
\centering \vspace{-0.3cm} \textbf{A.} 18\% & 
\centering \vspace{-0.3cm} \textbf{B.} 12\% & 
\centering \vspace{-0.3cm} \textbf{C.} 10\% & 
\centering \vspace{-0.3cm} \textbf{D.} 15\% \tabularnewline
\hline
\end{tabular}

\vspace{0.7cm}

\textbf{Question 10 :}
\vspace{5pt}

Un prix subit deux augmentations successives de 5\% puis de 12\%. Le taux d'évolution global, exprimé en pourcentage, est :
\vspace{0.3cm}

\renewcommand{\arraystretch}{3}
\begin{tabular}{|m{3.4cm}|m{3.4cm}|m{3.4cm}|m{3.4cm}|}
\hline
\centering \vspace{-0.3cm} \textbf{A.} 17\% & 
\centering \vspace{-0.3cm} \textbf{B.} 18\% & 
\centering \vspace{-0.3cm} \textbf{C.} 16,5\% & 
\centering \vspace{-0.3cm} \textbf{D.} 17,6\% \tabularnewline
\hline
\end{tabular}

\end{document}
